\documentclass[../../../include/open-logic-section]{subfiles}

\begin{document}

\olfileid{sth}{story}{approach}
\olsection{Pendekatan Kumulatif-Iteratif}

Iki panjelasan kang luwih jangkep sethithik babagan cara kita
merang himpunan dadi tataran:
\begin{quote}
	Himpunan kawangun ing \emph{tataran-tataran}. Kanggo saben
	tataran $S$, ana tataran tartamtu kang \emph{sadurunge} $S$.
	Ing tataran $S$, saben koleksi kang anggotane himpunan kang
	kawangun ing tataran sadurunge $S$ diwangun dadi sawijining
	himpunan. Ora ana himpunan liyane kajaba himpunan kang
	kawangun ing tataran-tataran. \citep[p.~323]{Shoenfield:AST}
\end{quote}
Iki gambaran cekak \emph{konsepsi kumulatif-iteratif
tumrap himpunan}. Konsepsi iki bakal dadi dhasar teori himpunan
formal kang kita aturake ing \olref[sth][][]{part}.

Saiki ayo dideleng luwih rinci. Kaya panjelasane Shoenfield,
ing saben tataran kita nggawe himpunan anyar saka himpunan
kang wis kasedhiya saka tataran sadurunge. Dadi ing gambaran
Shoenfield, tataran wiwitan, yaiku tataran $0$, ora nduweni
tataran \emph{sadurunge}; mula \emph{luwih-lebih} ora ana
himpunan saka tataran sadurunge kang kasedhiya.\footnote{Yagene
kita kudu nganggep \emph{ana} tataran kapisan? Delengen
cathetan ing \stagesord{} sajrone \olref[z][story]{sec}.}
Mula kita mung nggawe siji himpunan: himpunan tanpa
!!{element}s, yaiku $\emptyset$. Ing tataran $1$, persis
siji himpunan kasedhiya saka tataran sadurunge, mula mung
siji himpunan anyar, yaiku $\{\emptyset\}$. Ing tataran $2$,
loro himpunan kasedhiya saka tataran sadurunge, lan kita
nggawe rong himpunan anyar $\{\{\emptyset\}\}$ lan
$\{\emptyset, \{\emptyset\}\}$. Ing tataran $3$, papat
himpunan wis kasedhiya saka tataran sadurunge, mula kita
nggawe rolas himpunan anyar\ldots. Dadi gambaran
kumulatif-iteratif tumrap himpunan kaya ing ngisor iki
(angka nuduhake tataran):
\begin{center}
	\begin{tikzpicture}[scale=0.6]
	\tikzset{cut_here/.style={densely dotted}}
	\draw (-4,6) -- (-1, 0)-- (2,6); 
	\draw[cut_here] (-5,8)--(-4,6);
	\draw[cut_here] (2,6)--(3,8);
	\draw(-1.5, 1)--(-0.5, 1);
	\draw(-2, 2)--(0, 2);
	\draw(-2.5, 3)--(0.5,3);
	\draw(-3, 4)--(1, 4);
	\draw(-3.5, 5)--(1.5, 5);
	\draw(-4, 6)--(2, 6);
	\node[label] at (0, 0) {\small 0};
	\node[label] at (0.5, 1) {\small 1};
	\node[label] at (1, 2) {\small 2};
	\node[label] at (1.5, 3) {\small 3};
	\node[label] at (2, 4) {\small 4};
	\node[label] at (2.5, 5) {\small 5};
	\node[label] at (3, 6) {\small 6};
	\end{tikzpicture}
\end{center}
Mula, yagene kita kudu nampa gambaran iki?

Alesan siji yaiku gambaran iki apik lan gampang digarap.
Amarga gambaran kang paling cetha, yaiku Komprehensi Naif,
wis gagal, kita mbutuhake gambaran kang apik.

Nanging gambaran iki ora \emph{mung} apik. Kita nduweni
alesan kang becik kanggo percaya yen teori himpunan kang
adhedhasar gambaran iki bakal \emph{konsisten}. Iki alesane.

Miturut konsepsi kumulatif-iteratif, kita nggawe himpunan
ing tataran-tataran, lan !!{element}s-e kudu objek kang
\emph{wis} kasedhiya. Mula kanggo saben tataran~$S$,
kita bisa nggawe himpunan
\[
	R_S = \Setabs{x}{x \notin x \text{ lan $x$ wis kasedhiya sadurunge $S$}}
\]
Penalaran kang digunakake kanggo mbuktekake Paradoks Russell
saiki nuduhake yen $R_S$ dhewe durung kasedhiya sadurunge
tataran $S$. Iki dudu kontradiksi. Kajaba iku, yen kita
nampa konsepsi kumulatif-iteratif, kita kudune ora
\emph{ngarepake} bisa nggawe himpunan Russell dhewe. Amarga
himpunan iku bakal ngemot kabeh himpunan kang ora dadi
anggota awake dhewe lan kang ``bakal tau kasedhiya''.
Ringkese: miturut gambaran kumulatif-iteratif, kasunyatan
yen kita (kanthi bukti) ora bisa nggawe himpunan Russell
ora \emph{nggumunake}; iki pancen kang bakal kita
\emph{prakirakake}.

%Ing sawijine sisi, konsepsi kumulatif-iteratif tumrap himpunan menehi wangsulan marang Paradoks Russell kang rada padha karo wangsulane \emph{pendukung predikativitas}. Amarga pendukung predikativitas nganggep koleksi kabeh himpunan$_n$ kang ora dadi anggota awake dhewe minangka himpunan$_{n+1}$, lan himpunan$_n$ iki ora dadi anggota awake dhewe. (Malah akeh pendukung predikativitas nganggep pitakon apa sawijining himpunan dadi anggota awake dhewe minangka pitakon kang \emph{ora gramatikal}.)
%
%Nanging ana prabeda wigati antarane pendekatan kumulatif-iteratif lan pendekatan predikativis: pendekatan kumulatif-iteratif nganggep kabeh entitas \emph{sajinis}. Pendukung predikativitas mbokmenawa nduweni himpunan kosong$_0$ (yaiku himpunan$_0$ tanpa !!{element}s) lan himpunan kosong$_1$ (yaiku himpunan$_1$ tanpa !!{element}s), lan loro iki dadi \emph{entitas kang beda}. Nanging ing pendekatan kumulatif-iteratif mung ana siji himpunan kosong, $\emptyset$, kang kasedhiya ing saben tataran hierarki.

%Pancen, Aksioma Ekstensionalitas kang diandharake ing wiwitan bab iki bakal lumaku kanggo unsur-unsur ing hierarki iki (mula mung bisa ana \emph{siji} himpunan ``kosong''). Nanging konsepsi kumulatif-iteratif ora dimaksudake kanggo mbenerake Ekstensionalitas (delengen \cite{Boolos1971}). Sepisan maneh: Ekstensionalitas becike ditampa amarga kita ngrembug koleksi ekstensional.

\end{document}
